\uuid{eej0}
\titre{Rang d'une famille de vecteurs/Polynômes}
\auteur{hanani ; mzari}
\organisation{exo7}
\licence{CC BY-NC-SA 4.0}
\datecreate{2020-01-20}
\niveau{L1}
\module{Algèbre}
\chapitre{Espace vectoriel}
\sousChapitre{Base}
\difficulte{3}
\theme{Base et dimension, Rang d'une famille de vecteurs, Polynômes}
\exo7qcmid{lille-2.espaces-vectoriels.038}

\contenu{
\begin{question}
\qclassification{106.04}
\qtags{motcle=Rang d'une famille de vecteurs/Polynômes}
Dans $\Rr_3[X]$, l'ensemble des polynômes à coefficients réels de degré $\le 3$, on considère les polynômes $P_1= X^3+1, P_2= P'_1 $ (la dérivée de $P_1$) et  $ P_3 = P''_1$ (la dérivée seconde  de $P_1$). Quelles sont les assertions vraies ?
\begin{answers}
\bad{Le rang de la famille $\{P_1, P_3 \}$ est $3$}.
\bad{$\{P_1,P_2,P_3 \}$ est une famille génératrice de $\Rr_3[X]$}.
\good{$\{P_1,P_2,P_3 \}$ est une famille libre de $\Rr_3[X]$}.
\good{Le rang de la famille  $\{P_1, P_2,P_3 \}$ est $3$}.
\end{answers}
\begin{explanations}  On vérifie que $\{P_1,P_2,P_3 \}$ est une famille libre de $\Rr_3[X]$ (ce sont des polynômes de degrés
distincts). Par contre, elle n'est pas génératrice de $\Rr_3[X]$, puisque la dimension de cet espace est $4$.
\vskip2mm
Le rang d'une famille de vecteurs est la dimension du sous-espace vectoriel engendré par ces vecteurs. Autrement dit, c'est le nombre maximum de vecteurs linéairement indépendants de cette famille.
\end{explanations}
\end{question}
}
